\chapter{Comprehensive Questionnaire Validation}
\label{ch:comprehensive}

This chapter clarifies the relationship between a global satisfiability check of the questionnaire and the per-item implication checks used in the implementation. We formalize both notions and prove what is (and is not) equivalent, with practical consequences for validation.

\section{Setup}

Recall from \Cref{def:domain-constraint} that each outcome variable $S_i$ has an associated domain constraint $D_i(S_i)$. Let $B := \bigwedge_{i=1}^n D_i(S_i)$ denote the conjunction of all domain constraints in the questionnaire. For each item $I_i$, let $P_i$ be its precondition and $Q_i$ its postcondition, expressed as Boolean formulas over the answers $\mathbf{S} = (S_1,\dots,S_n)$.

\begin{definition}[Global Satisfiability Formula]
\label{def:global-formula}
The \textbf{global satisfiability formula} for a questionnaire is:
\[
  \mathcal{F} \;:=\; B \;\wedge\; \bigwedge_{i=1}^n \bigl(P_i \Rightarrow Q_i\bigr)
\]
We say the questionnaire passes the global satisfiability check if $\mathrm{SAT}(\mathcal{F})$ holds.
\end{definition}

\begin{definition}[Witness Formula]
\label{def:witness-formula}
For an item $I_i$, the \textbf{witness formula} is:
\[
  W_i := B \wedge P_i \wedge \lnot Q_i
\]
If $\mathrm{SAT}(W_i)$, then there exists an assignment satisfying the domain constraints where item $I_i$ is reached (precondition $P_i$ holds) but the postcondition $Q_i$ is violated. Such an assignment demonstrates that the item can fail validation.
\end{definition}

The key diagnostic question is whether $W_i$ is satisfiable: $\mathrm{UNSAT}(W_i)$ means item $I_i$ is \emph{locally valid} (whenever it is reached, its postcondition must hold), while $\mathrm{SAT}(W_i)$ indicates a potential violation.

\section{Key Relations}

The following theorems capture the exact relationship between the global check and the per-item checks.

\begin{theorem}[Soundness of Per-Item Validity]
\label{thm:soundness}
If $\mathrm{SAT}(B)$ and for every $i \in \{1, \ldots, n\}$ we have $\mathrm{UNSAT}(W_i)$, then $\mathrm{SAT}(\mathcal{F})$.
\end{theorem}

\begin{proof}
Assume $\mathrm{SAT}(B)$ and $\mathrm{UNSAT}(W_i)$ for all $i \in \{1, \ldots, n\}$, where $W_i = B \wedge P_i \wedge \lnot Q_i$.

\textbf{Step 1: Transform unsatisfiability to entailment for each item.}

For each $i$, we have $\mathrm{UNSAT}(B \wedge P_i \wedge \lnot Q_i)$.

By definition of unsatisfiability:
\[
  \nexists \mathbf{a} : \mathbf{a} \models (B \wedge P_i \wedge \lnot Q_i)
\]

Equivalently, by the quantifier negation rule (i.e., $\neg \exists x \, P(x) \equiv \forall x \, \neg P(x)$):
\[
  \forall \mathbf{a} : \mathbf{a} \not\models (B \wedge P_i \wedge \lnot Q_i)
\]

By the definition of satisfaction for conjunction (an assignment fails to satisfy a conjunction if it fails to satisfy at least one conjunct), this means for all assignments $\mathbf{a}$:
\[
  \mathbf{a} \not\models B \text{ or } \mathbf{a} \not\models P_i \text{ or } \mathbf{a} \not\models \lnot Q_i
\]

Note that $\mathbf{a} \not\models \lnot Q_i$ is equivalent to $\mathbf{a} \models Q_i$ (by double negation).

Rewriting the disjunction using the logical equivalence $(\neg A \vee \neg B \vee C) \equiv ((A \wedge B) \Rightarrow C)$, we get: for all assignments $\mathbf{a}$:
\[
  \text{if } \mathbf{a} \models B \wedge P_i, \text{ then } \mathbf{a} \models Q_i
\]

by definition of $\models$, this is equivalent to:
\[
  B \wedge P_i \models Q_i
\]

By the semantic deduction theorem (moving $P_i$ to the right), this is the definition of:
\[
  B \models (P_i \Rightarrow Q_i)
\]

\textbf{Step 2: Combine entailments for all items.}

We have established that $B \models (P_i \Rightarrow Q_i)$ for each $i \in \{1, \ldots, n\}$.

By the monotonicity of entailment: if $B \models \phi_i$ for each $i$, then $B \models \bigwedge_i \phi_i$.

Therefore:
\[
  B \models \bigwedge_{i=1}^n (P_i \Rightarrow Q_i)
\]

\textbf{Step 3: Construct a model for $\mathcal{F}$.}

By hypothesis, $\mathrm{SAT}(B)$, so there exists an assignment $\mathbf{a}$ such that $\mathbf{a} \models B$.

From Step 2, we have $B \models \bigwedge_{i=1}^n (P_i \Rightarrow Q_i)$.

By the definition of semantic entailment, every model of $B$ is also a model of $\bigwedge_{i=1}^n (P_i \Rightarrow Q_i)$.

Since $\mathbf{a} \models B$, we have $\mathbf{a} \models \bigwedge_{i=1}^n (P_i \Rightarrow Q_i)$.

Therefore:
\[
  \mathbf{a} \models B \wedge \bigwedge_{i=1}^n (P_i \Rightarrow Q_i) = \mathcal{F}
\]

Hence, $\mathrm{SAT}(\mathcal{F})$.
\end{proof}

\begin{theorem}[Local Invalidity Does Not Imply Global Unsatisfiability]
\label{thm:local-not-global}
If $\mathrm{SAT}(W_i)$ for some $i$, it does not follow that $\mathrm{UNSAT}(\mathcal{F})$.
\end{theorem}

The witness $W_i$ only shows that there \emph{exists} a violating model of $B$; the global check asks whether there \emph{exists} a (possibly different) model that satisfies all implications simultaneously.

\begin{mdframed}
\noindent\textbf{Example.} Consider a two-item questionnaire:
\begin{itemize}
    \item $I_1$: ``What is your age?'' with $S_1 \in [0, 120]$, $P_1 = \mathit{true}$, $Q_1 = \mathit{true}$
    \item $I_2$: ``Years of driving experience?'' with $S_2 \in [0, 100]$, $P_2 = (S_1 \geq 16)$, $Q_2 = (S_2 \leq S_1 - 16)$
\end{itemize}

The postcondition $Q_2$ ensures driving experience cannot exceed years since turning 16. Consider the witness formula for $I_2$:
\[
W_2 = B \wedge P_2 \wedge \lnot Q_2 = B \wedge (S_1 \geq 16) \wedge (S_2 > S_1 - 16)
\]

This is satisfiable: $S_1 = 20$, $S_2 = 10$ satisfies $W_2$ (claiming 10 years of experience at age 20, which violates $Q_2$ since $10 > 20 - 16 = 4$). So $\mathrm{SAT}(W_2)$—the per-item check reports item $I_2$ as \emph{locally invalid}.

However, the global formula $\mathcal{F}$ is still satisfiable. Consider $S_1 = 30$, $S_2 = 5$:
\begin{itemize}
    \item $P_2 = (30 \geq 16) = \mathit{true}$, and $Q_2 = (5 \leq 30 - 16) = (5 \leq 14) = \mathit{true}$
    \item Thus $(P_2 \Rightarrow Q_2) = \mathit{true}$
\end{itemize}

This assignment satisfies $\mathcal{F}$, so $\mathrm{SAT}(\mathcal{F})$ despite $\mathrm{SAT}(W_2)$. The local invalidity witness shows that \emph{some} respondents could violate the postcondition, but the questionnaire still has valid completions.
\end{mdframed}

\begin{theorem}[Accumulated Constraints Can Cause Global Unsatisfiability]
\label{thm:accumulated-unsat}
There exist questionnaires where no item is INFEASIBLE (each $\mathrm{SAT}(B \wedge P_i \wedge Q_i)$), yet the global formula $\mathcal{F}$ is unsatisfiable.
\end{theorem}

Per-item checks examine each postcondition independently against the domain constraints $B$. However, during questionnaire execution, items are evaluated in dependency order, and postconditions from earlier items must be satisfied before proceeding to later items. These accumulated constraints are not visible to per-item checks.

\begin{mdframed}
\noindent\textbf{Example.} Consider a two-item questionnaire where $I_2$ depends on $S_1$:
\begin{itemize}
    \item $I_1$: ``Rate your experience (1--100)'' with $S_1 \in [1, 100]$, $P_1 = \mathit{true}$, $Q_1 = (S_1 > 50)$
    \item $I_2$: ``Confirm your rating'' with $S_2 \in \{0, 1\}$, $P_2 = \mathit{true}$, $Q_2 = (S_1 < 30)$
\end{itemize}

Since $Q_2$ references $S_1$, the dependency graph requires $I_1 \to I_2$ ordering. The respondent must satisfy $Q_1$ (i.e., $S_1 > 50$) before proceeding to $I_2$.

\textbf{Per-item analysis:}
\begin{itemize}
    \item $W_1 = B \wedge \mathit{true} \wedge (S_1 \leq 50)$ is SAT (e.g., $S_1 = 25$) $\Rightarrow$ CONSTRAINING
    \item $W_2 = B \wedge \mathit{true} \wedge (S_1 \geq 30)$ is SAT (e.g., $S_1 = 50$) $\Rightarrow$ CONSTRAINING
    \item Neither item is INFEASIBLE: $\mathrm{SAT}(B \wedge P_1 \wedge Q_1)$ and $\mathrm{SAT}(B \wedge P_2 \wedge Q_2)$
\end{itemize}

\textbf{Global analysis:}
\[
\mathcal{F} = B \wedge (\mathit{true} \Rightarrow S_1 > 50) \wedge (\mathit{true} \Rightarrow S_1 < 30) = B \wedge (S_1 > 50) \wedge (S_1 < 30)
\]

This is unsatisfiable—no value of $S_1$ can simultaneously satisfy $S_1 > 50$ and $S_1 < 30$. Therefore $\mathrm{UNSAT}(\mathcal{F})$.

The per-item checks miss this conflict because they evaluate each postcondition against all possible values of $S_1$, not accounting for the fact that $Q_1$ must already be satisfied when $I_2$ is reached.
\end{mdframed}

\section{Path-Based Validation}

Path-based validation addresses a limitation of the global formula $\mathcal{F}$: the use of implications $(P_i \Rightarrow Q_i)$ that are vacuously true when preconditions are false. This section formalizes the path-based approach and proves its relationship to global validation.

\begin{definition}[Execution Path]
\label{def:execution-path}
An \textbf{execution path} $\pi$ through a questionnaire is a sequence of items $(I_{k_1}, I_{k_2}, \ldots, I_{k_m})$ respecting the dependency order, where each item's precondition is satisfied given the accumulated constraints from preceding items.
\end{definition}

\begin{definition}[Path Constraint Formula]
\label{def:path-formula}
For an execution path $\pi = (I_{k_1}, \ldots, I_{k_m})$, the \textbf{path constraint formula} is:
\[
  \Phi_\pi \;:=\; B \;\wedge\; \bigwedge_{j=1}^{m} \bigl(P_{k_j} \wedge Q_{k_j}\bigr)
\]
The path is \textbf{feasible} if $\mathrm{SAT}(\Phi_\pi)$.
\end{definition}

Note the key difference from the global formula: path constraints use conjunctions $(P_i \wedge Q_i)$ for items on the path, not implications $(P_i \Rightarrow Q_i)$. This ensures that preconditions are actually satisfiable, not just vacuously handled.

\begin{theorem}[Global Validity is Necessary for Path Validity]
\label{thm:global-necessary}
If $\mathrm{UNSAT}(\mathcal{F})$, then for every execution path $\pi$, we have $\mathrm{UNSAT}(\Phi_\pi)$.
\end{theorem}

\begin{proof}
For any path $\pi = (I_{k_1}, \ldots, I_{k_m})$, if $\mathrm{SAT}(\Phi_\pi)$, then there exists an assignment $\mathbf{a}$ satisfying $B \wedge \bigwedge_{j=1}^m (P_{k_j} \wedge Q_{k_j})$.

For items on the path: $\mathbf{a} \models P_{k_j}$ and $\mathbf{a} \models Q_{k_j}$, so $\mathbf{a} \models (P_{k_j} \Rightarrow Q_{k_j})$.

For items not on the path: their preconditions are false under $\mathbf{a}$ (otherwise they would be on the path), so $(P_i \Rightarrow Q_i)$ is vacuously true.

Therefore $\mathbf{a} \models \mathcal{F}$, contradicting $\mathrm{UNSAT}(\mathcal{F})$.
\end{proof}

\begin{theorem}[Global Validity Does Not Guarantee All Items Reachable]
\label{thm:global-not-sufficient}
There exist questionnaires where $\mathrm{SAT}(\mathcal{F})$ but some item $I_i$ with $P_i \neq \mathit{false}$ is unreachable on every feasible path.
\end{theorem}

\begin{mdframed}
\noindent\textbf{Example.} Consider:
\begin{itemize}
    \item $I_1$: $S_1 \in [1, 100]$, $P_1 = \mathit{true}$, $Q_1 = (S_1 > 80)$
    \item $I_2$: $S_2 \in \{0, 1\}$, $P_2 = (S_1 < 50)$, $Q_2 = \mathit{true}$
\end{itemize}

The dependency graph requires $I_1 \to I_2$ (since $P_2$ references $S_1$).

\textbf{Per-item analysis:}
\begin{itemize}
    \item $P_2 = (S_1 < 50)$ is CONDITIONAL: $\mathrm{SAT}(B \wedge S_1 < 50)$ with $S_1 = 25$
    \item $Q_1 = (S_1 > 80)$ is CONSTRAINING: $\mathrm{SAT}(W_1)$ with $S_1 = 25$
\end{itemize}

\textbf{Global analysis:}
\[
\mathcal{F} = B \wedge (\mathit{true} \Rightarrow S_1 > 80) \wedge ((S_1 < 50) \Rightarrow \mathit{true}) \equiv B \wedge (S_1 > 80)
\]
This is satisfiable (e.g., $S_1 = 90$), so $\mathrm{SAT}(\mathcal{F})$.

\textbf{Path-based analysis:}
Any path including $I_2$ requires:
\[
\Phi_\pi = B \wedge (P_1 \wedge Q_1) \wedge (P_2 \wedge Q_2) \equiv B \wedge (S_1 > 80) \wedge (S_1 < 50)
\]
This is unsatisfiable—no value of $S_1$ can satisfy both $S_1 > 80$ and $S_1 < 50$.

Therefore, $I_2$ is classified as CONDITIONAL by per-item analysis (reachable when $S_1 < 50$), the global check passes, but $I_2$ is actually unreachable on any valid execution path. This is \textbf{dead code} that path-based validation detects.
\end{mdframed}

\paragraph{Practical Path Enumeration.}

In practice, path enumeration leverages the dependency graph structure constructed in \Cref{sec:dep-graph-cycle}. The topological ordering computed by Kahn's algorithm (\Cref{alg:kahn}) determines the evaluation order: items with no dependencies can be traversed in any order, while dependent items must follow their predecessors. The path-based check:

\begin{enumerate}
    \item Uses the dependency graph from \Cref{sec:dep-graph-cycle} to identify predecessor relationships
    \item For each item $I_i$ with a CONDITIONAL precondition, constructs the accumulated constraint formula from all predecessor items in the topological order
    \item Checks if the item remains reachable under accumulated constraints
\end{enumerate}

\begin{definition}[Accumulated Reachability]
\label{def:accumulated-reachability}
For item $I_i$, let $\mathrm{Pred}(i) = \{j : I_j \prec I_i\}$ be the set of predecessor items, where $I_j \prec I_i$ denotes that $I_j$ precedes $I_i$ in the topological order. The \textbf{accumulated constraint formula} is:
\[
  \mathcal{A}_i \;:=\; B \;\wedge\; \bigwedge_{j \in \mathrm{Pred}(i)} (P_j \Rightarrow Q_j)
\]
Item $I_i$ is \textbf{accumulated-reachable} if $\mathrm{SAT}(\mathcal{A}_i \wedge P_i)$.
\end{definition}

\begin{mdframed}
\noindent\textbf{Example.} Consider a three-item questionnaire:
\begin{itemize}
    \item $I_1$: ``Annual income'' with $S_1 \in [0, 1000000]$, $P_1 = \mathit{true}$, $Q_1 = (S_1 \geq 50000)$
    \item $I_2$: ``Tax bracket'' with $S_2 \in \{1, 2, 3\}$, $P_2 = \mathit{true}$, $Q_2 = \mathit{true}$
    \item $I_3$: ``Low-income assistance'' with $S_3 \in \{0, 1\}$, $P_3 = (S_1 < 30000)$, $Q_3 = \mathit{true}$
\end{itemize}

The dependency graph has $I_1 \to I_3$ (since $P_3$ references $S_1$), while $I_2$ is independent.

\textbf{Per-item analysis:}
\begin{itemize}
    \item $P_3 = (S_1 < 30000)$ is CONDITIONAL: $\mathrm{SAT}(B \wedge S_1 < 30000)$ holds
\end{itemize}

\textbf{Accumulated reachability for $I_3$:}

The predecessor set is $\mathrm{Pred}(3) = \{1\}$ (only $I_1$ precedes $I_3$ in the dependency order; $I_2$ is independent and not a predecessor).

\[
\mathcal{A}_3 = B \wedge (P_1 \Rightarrow Q_1) = B \wedge (\mathit{true} \Rightarrow S_1 \geq 50000) = B \wedge (S_1 \geq 50000)
\]

Now check accumulated reachability:
\[
\mathcal{A}_3 \wedge P_3 = B \wedge (S_1 \geq 50000) \wedge (S_1 < 30000)
\]

This is unsatisfiable—no income value can be both $\geq 50000$ and $< 30000$. Therefore $I_3$ is \textbf{not accumulated-reachable}.

The questionnaire designer intended $I_3$ (low-income assistance) to appear for respondents with income below \$30,000, but $Q_1$ requires income $\geq$ \$50,000 to proceed. This is dead code that accumulated reachability analysis detects.
\end{mdframed}

If $\mathrm{UNSAT}(\mathcal{A}_i \wedge P_i)$ for a CONDITIONAL item, the item is dead code—its precondition can never be satisfied given the constraints imposed by predecessor items.

\section{Validation Hierarchy}

Questionnaire validation can be performed at three levels of increasing thoroughness and computational cost:

\begin{enumerate}
    \item \textbf{Per-item validation} (\Cref{ch:definitions}): For each item $I_i$, check independently:
    \begin{itemize}
        \item \emph{Reachability}: Classify $P_i$ as ALWAYS, NEVER, or CONDITIONAL using $\mathrm{SAT}(B \wedge P_i)$
        \item \emph{Validity}: Classify $Q_i$ as TAUTOLOGICAL, CONSTRAINING, or INFEASIBLE using the witness formula $W_i$
    \end{itemize}
    This is the fastest level—each item can be checked independently and in parallel.

    \item \textbf{Global validation}: Check $\mathrm{SAT}(\mathcal{F})$ where $\mathcal{F} = B \wedge \bigwedge_i (P_i \Rightarrow Q_i)$.
    This verifies that at least one valid questionnaire completion exists, but uses implications that are vacuously true when preconditions are false.

    \item \textbf{Path-based validation}: Enumerate all execution paths through the questionnaire and check each path's accumulated constraints. This is the most thorough level, detecting items that become unreachable due to accumulated postconditions (dead code).
\end{enumerate}

\paragraph{Relationships Between Levels.}

The three validation levels form a hierarchy with precise logical relationships:

\begin{itemize}
    \item \textbf{Per-item passes $\Rightarrow$ Global passes}: By \Cref{thm:soundness}, if all $W_i$ are UNSAT (all postconditions TAUTOLOGICAL), then $\mathrm{SAT}(\mathcal{F})$ is guaranteed. No further checks needed.

    \item \textbf{Global fails $\Rightarrow$ Path-based fails}: If $\mathrm{UNSAT}(\mathcal{F})$, then no execution path can be valid. The global formula is a necessary condition for questionnaire validity.

    \item \textbf{Global passes $\not\Rightarrow$ All paths valid}: The global check can pass while some items are unreachable, because implications $(P_i \Rightarrow Q_i)$ are vacuously true when $P_i$ is false.
\end{itemize}

\paragraph{When Each Level Suffices.}

\begin{itemize}
    \item \textbf{Per-item validation suffices} when all postconditions are TAUTOLOGICAL. The questionnaire has no constraining logic, so any answer combination is valid.

    \item \textbf{Global validation suffices} when you only need to verify that \emph{some} valid completion exists. Dead code (unreachable items) is acceptable.

    \item \textbf{Path-based validation is needed} when you require that \emph{every} conditional item is actually reachable on at least one valid path. This detects design errors where accumulated postconditions make items unreachable.
\end{itemize}

\paragraph{Summary of Validation Results.}

\begin{itemize}
    \item \textbf{All $W_i$ UNSAT}: All postconditions are TAUTOLOGICAL. By \Cref{thm:soundness}, global and path-based validity are guaranteed—no further checks needed.

    \item \textbf{For some $i$, $\mathrm{SAT}(W_i)$, $\mathrm{UNSAT}(\mathcal{F})$}: Accumulated CONSTRAINING postconditions conflict (\Cref{thm:accumulated-unsat}). By \Cref{thm:global-necessary}, no execution path is valid. This is a design error.

    \item \textbf{For some $i$, $\mathrm{SAT}(W_i)$, $\mathrm{SAT}(\mathcal{F})$, all items accumulated-reachable}: The questionnaire is fully valid—completable with no dead code.

    \item \textbf{For some $i$, $\mathrm{SAT}(W_i)$, $\mathrm{SAT}(\mathcal{F})$, some items not accumulated-reachable}: The questionnaire is completable (\Cref{thm:local-not-global}), but contains dead code (\Cref{thm:global-not-sufficient}). Report as warning.

    \item \textbf{Some item INFEASIBLE}: The postcondition cannot be satisfied when the item is reached—design error requiring fix.

    \item \textbf{Some item NEVER reachable (per-item)}: Precondition is unsatisfiable even without accumulated constraints—design error.
\end{itemize}

